% Generated draft: requires mathematical and source review.
\documentclass[11pt]{article}
\usepackage[margin=25mm]{geometry}
\usepackage{fontspec}
\setmainfont{Latin Modern Roman}
\usepackage{amsmath,amssymb,mathtools,unicode-math}
\setmathfont{Latin Modern Math}
\usepackage{enumitem,needspace,xurl,hyperref}
\setlength{\parindent}{0pt}
\setlength{\parskip}{5pt}
\setlength{\emergencystretch}{3em}
\newcommand{\N}{\mathbb{N}}
\newcommand{\Z}{\mathbb{Z}}
\newcommand{\Q}{\mathbb{Q}}
\newcommand{\R}{\mathbb{R}}
\newcommand{\C}{\mathbb{C}}
\newcommand{\F}{\mathbb{F}}
\newcommand{\A}{\mathbb{A}}
\newcommand{\PP}{\mathbb{P}}
\newcommand{\E}{\mathbb{E}}
\newcommand{\eps}{\varepsilon}
\newcommand{\dd}{\,\mathrm{d}}
\begin{document}
\section*{UnsolvedMath ID: 20000459}
\textbf{Source problem code:} catalogue entry
\subsection{Definitions and mathematical statement}
For a positive integer \(d\), let
\[
F_d=\{[x_1:\cdots:x_6]\in\PP^5_\C:x_1^d+\cdots+x_6^d=0\}.
\]
This is the intended Fermat fourfold interpretation of the damaged exponent typography; the supplied record marks the reconstruction unverified. Heath-Brown's Question 26 asks to describe the parameter spaces of curves of low degree on \(F_d\), with the curves taken over \(\C\). The source gives no numerical meaning of ``low degree'' and no choice of parameter scheme or restrictions on reducibility.

For a precise integral-curve part of that question, let \(C\subset F_d\) be a reduced irreducible projective curve and put \(\delta=\deg_{\mathcal O(1)}C\). For each Hilbert polynomial \(P(t)=\delta t+1-p_a\), one may ask to describe the integral-curve locus in \(\operatorname{Hilb}^{P}(F_d)\), where \(p_a\) is the arithmetic genus. A full answer to the original wording may require other curves and scheme structures too; those have not been silently excluded from its scope. The workshop remark says lines are known and expresses the hope of ``none or very few'' further curves, without explaining what range or exceptional families this refers to.

The attached aid proposes the following partial classification. Given a partition of \(\{1,\ldots,6\}\) into three unordered pairs \(\{a_i,b_i\}\), choose an ordering within each pair and numbers \(\zeta_i\in\C\) with \(\zeta_i^d=-1\). The equations
\[
x_{a_i}=\zeta_i x_{b_i}\quad(i=1,2,3)
\]
define a standard coordinate-pairing plane \(\Pi\subset F_d\). In the range in question there are \(15d^3\) such planes. The aid claims that, for every integral curve of degree \(\delta\geq2\),
\[
d>10\delta-6\quad\Longrightarrow\quad
C\text{ lies in a unique standard coordinate-pairing plane}.
\]
It describes the resulting set of closed points of the integral degree-\(\delta\) parameter locus as the disjoint union, over these planes, of the sets of irreducible homogeneous degree-\(\delta\) plane forms up to nonzero scalar. For each plane this is the irreducible-form open subset of
\(\PP(\C[u,v,w]_\delta)\). Such curves have arithmetic genus \((\delta-1)(\delta-2)/2\). Crucially, the aid makes no assertion identifying Hilbert-scheme nilpotents, embedded infinitesimal directions, or full parameter-scheme structures with those plane-form loci.

The same aid reports that every nonlinear integral rational curve is contained in a standard plane for \(d\geq24\), and that a curve of genus \(g\geq1\) is so contained if \(d>20g+4\). Whether this latter \(g\) means geometric or arithmetic genus is not specified in the summary and must be checked. For conics, \(d>10\cdot2-6\) is \(d\geq15\); the aid says that Reznick's published identity gives a nonstandard smooth conic on \(F_{14}\), making that cutoff sharp. Neither the identity nor the conic is supplied. All of these classifications are reported unreviewed research-aid claims, not independently established results of this draft.
\subsection{Short English statement}
Describe the low-degree complex curves on the Fermat fourfold and their parameter spaces. The attached aid proposes that sufficiently low-degree integral curves lie in unique standard coordinate-pairing planes, while explicitly limiting its parameter description to closed points and leaving Hilbert-scheme infinitesimal structure untreated.
\subsection{Sources}
\begin{itemize}
\item Heath-Brown, Question 26, AIM workshop \emph{Rational and integral points on higher dimensional varieties}: \url{https://aimath.org/WWN/qptsurface2/qptsurface2.pdf}. The supplied verification locates Question 26 on PDF p. 53 but does not append the promised restored statement.
\item Local registry \texttt{data/unsolvedmath/problems.json}, permanent ID \(20000459\), exact source code \texttt{catalogue entry}; canonical \texttt{aim-algebraic-number-theory-notes.json}, notes[110].
\item The attached aid invokes Salberger's diagonal-curve inequality and Reznick's published identity without full bibliographic references. Attempt 1 is explicitly machine-generated, not peer reviewed, and requires independent verification before citation.
\end{itemize}
\end{document}
